\subsection{STATIONARY SEQUENCES}

DEFINITION 5. A sequence \((x_{n})_{n\in \mathbb{N}}\) of elements of a set \(\mathbf{E}\) is said to be stationary if there exists an integer \(m\) such that \(x_{n} = x_{m}\) for all integers \(n\geqslant m\) .

PROPOSITION 6. Let E be an ordered set. Then the following statements are equivalent:

\begin{enumerate}
\def\labelenumi{(\alph{enumi})}
\item
  Every non-empty subset of \(\mathbf{E}\) has a maximal element.
\item
  Every increasing sequence \((x_{n})\) of elements of E is stationary.
\end{enumerate}

We show first that (a) implies (b). Let X be the set of elements of the sequence \((x_{n})\) , and let \(x_{m}\) be a maximal element of X. If \(n \geqslant m\) , we have by hypothesis \(x_{n} \geqslant x_{m}\) , and therefore \(x_{n} = x_{m}\) by the maximality of \(x_{m}\) . Conversely, suppose that there exists a non-empty subset A of E which has no maximal element. For each \(x \in A\) , let \(T_{x}\) be the set of all \(y \in A\) such that y \textgreater{} x. By assumption, \(T_{x} \neq \emptyset\) for all \(x \in A\) ; hence there exists a mapping f of A into A such that \(f(x) > x\) for all \(x \in A\) (Chapter II, §5, no. 4, Proposition 6). If \(a \in A\) , then the sequence \((x_{n})_{n \in \mathbb{N}}\) defined inductively by the conditions \(x_{0} = a\) , \(x_{n+1} = f(x_{n})\) is evidently increasing and not stationary.

COROLLARY 1. A totally ordered set E is well-ordered if and only if every decreasing sequence of elements of E is stationary.

For to say that E is well-ordered is equivalent to saying that every non-empty subset of E has a minimal element ( \(\S\) 1, no. 10, Proposition 10), and the assertion therefore follows from Proposition 6.

COROLLARY 2. Every increasing sequence of elements of a finite ordered set is stationary.

For every finite ordered set has a maximal element (§ 4, no. 4, Proposition 3, Corollary 2).

¶ An ordered set E which satisfies the equivalent conditions of Proposition 6 is sometimes said to be Noetherian.

PROPOSITION 7 (``Principle of Noetherian induction''). Let E be a Noetherian set, and let F be a subset of E with the following property: if \(a \in E\) is such that the relation \(x > a\) implies \(x \in F\) , then \(a \in F\) . Under these conditions, \(F = E\) .

Indeed, suppose \(\mathbf{E} \neq \mathbf{F}\) ; then \(\mathbf{E} - \mathbf{F}\) has a maximal element \(b\) . By definition we have \(x \in \mathbf{F}\) for all \(x > b\) ; but this implies \(b \in \mathbf{F}\) , which is absurd.

